\section{Related Work}

\subsection{3D Visuomotor Policy Learning}

Imitation learning acquires manipulation skills from visual and proprioceptive demonstrations through behavioral cloning~\cite{florence2022implicit,zhao2023learning} or diffusion-based action generation~\cite{chi2025diffusion}. Recent policies make spatial structure explicit using point clouds~\cite{ze20243d,ke20243d}, voxel grids~\cite{shridhar2023perceiver}, or geometrically related multi-view features~\cite{goyal2023rvt}. Point clouds are particularly attractive because they preserve metric coordinates through a compact policy interface. However, raw point clouds do not identify functional parts, and RGB-D samples vary with viewpoint, occlusion, segmentation, and resampling. Existing policies typically absorb this variability through augmentation or stronger encoders. Our work instead keeps the point-cloud policy and action objective fixed while improving the object representation supplied to it.

\subsection{Semantic 3D Representations for Manipulation}

Semantic 3D representations supplement geometry with functional cues. One family of methods lifts pretrained 2D descriptors such as DINOv2~\cite{oquab2023dinov2} into 3D, including Distilled Feature Fields~\cite{shen2023distilled}, G3Flow~\cite{chen2025g3flow}, and GenDP~\cite{wang2024gendp}. Another family learns affordances, action trajectories, or dense part-aware features directly from 3D data~\cite{mo2021where2act,wu2021vat,chen2026learning,xu2026action}. Large-scale point-cloud encoders~\cite{xie2020pointcontrast,zhang2026utonia} and fine-grained part datasets such as PartNext~\cite{wang2026partnext} further support transferable geometric and part understanding. These approaches show that semantic features can expose task-relevant structure that geometry alone may not reveal. Nevertheless, descriptors evaluated only on currently observed points remain coupled to their visibility and sampling. We instead use part supervision to learn an object-conditioned function from which semantic embeddings are queried for policy input.


\subsection{Queryable Object-Centric Fields}
Object-centric correspondences provide a structured basis for category-level manipulation. kPAM represents manipulation-relevant keypoint affordances~\cite{manuelli2019kpam}, TAX-Pose estimates task-specific cross-object transformations~\cite{pan2023tax}, and Neural Descriptor Fields provide continuous SE(3)-equivariant descriptors~\cite{simeonov2022neural}. Dense feature fields, including D3Fields~\cite{wang2023d}, Distilled Feature Fields~\cite{shen2023distilled}, GenDP~\cite{wang2024gendp}, and G3Flow~\cite{chen2025g3flow}, further demonstrate the utility of queryable 3D semantics for learned manipulation. Our method shares the continuous-query formulation but targets the coupling between semantic descriptors and noisy observed samples. The observed object points condition the field, while the policy input is produced at resampled query locations. Category-level part objectives organize these queried embeddings around functional regions. We analyze the resulting representation and evaluate it in simulated and real-world manipulation without changing the downstream policy formulation.
